\scanpage{117}
\mspage{108}
% archive sheet number: 117 (lower left)
% L1-3
\uline{Théorème} 1) Tout support constructible $S$ s'écrit sous la forme $\operatorname{supp}^{\circ}\Phi$, où $\Phi$ est une préfigure non morcelée\textsuperscript{*};
% slanted
\margin{* Je dis plutôt : ``\unsure{irredondante|irrédondante}''}
% L4-6
\deleted{\unsure{2)} Cette écriture est unique, et on a}
et pour une telle $\Phi$, on a
\[
\operatorname{supp}^{\circ}(\Phi)=\operatorname{Omb}^{\circ}(\unsure{\Psi|\Phi})=\bigcup_{X\in\Phi}\operatorname{Omb}^{\circ}(\deleted{\unsure{\Phi}}X)
\]
(\uuline{réunion} \uuline{disjointe}).
% L7-11
Les éléments de $\inserted{\Phi_1=}\Phi\cap\unsure{M|\mathcal{M}}_1$ sont les éléments de $S$ qui sont \deleted{maximaux} $1^{\circ})$ de dim 1 et $2^{\circ})$ maximaux dans $S$ pour $\overset{\circ}{\ll}$. \inserted{(ou pour $\ll$, c'est pareil)} Les éléments de $\Phi_0=\Phi\cap M_0$ sont \deleted{les éléments de $S\cap M_0$ qui sont \unsure{suivants} : \unsure{où|ou} bien ils sont} \inserted{caractérisés par}
\[
\Phi_0=\left(\bigcup_{X\in\Phi_1}\unsure{\partial|\delta}X\right)\cap\deleted{\unsure{\Phi}}S\cup\Phi_{0\,\mathrm{is}}\quad\deleted{\unsure{\phi}}
\]
% L12-16
(réunion disjointe), où \deleted{$\Phi_{0\,\mathrm{is}}$ est}
\[
\begin{aligned}
\Phi_{0\,\mathrm{is}}&=\text{ens des sommets isolés de }\Phi\\
&=\text{ens. des éléments de }S\cap M_0\text{ qui}\\
&\qquad\text{sont maximaux dans }S\text{ pour }\ll
\end{aligned}
\]
\deleted{\unsure{Ce qui}}

% L17-22
\deleted{\unsure{Ainsi}} En particulier, $\Phi$ est bien \unsure{déterminé|déterminée} par la connaissance de $S$. On dit que $S$ est \uline{fermé} si $\Phi$ est fermée, i.e. est une figure --- \unsure{ou.|ou} si $S$ est le support d'une figure.
